% 巨行星（综述）
% license CCBYSA3
% type Wiki

本文根据 CC-BY-SA 协议转载翻译自维基百科\href{https://en.wikipedia.org/wiki/Giant_planet}{相关文章}。

\begin{figure}[ht]
\centering
\includegraphics[width=8cm]{./figures/6d40f7553e942e4f.png}
\caption{太阳系的四颗巨行星如下：（上）木星与土星（气体巨行星）（下）天王星与海王星（冰巨行星）它们按照与太阳的距离顺序排列，并以真实色彩呈现。尺寸不按比例绘制。} \label{fig_Giant_1}
\end{figure}
巨行星（giant planet），有时也被称为类木行星（jovian planet）（Jove 是罗马神话中木星〔Jupiter〕的另一名称），是一类比地球大得多、类型多样的行星。巨行星通常主要由低沸点物质（挥发物）构成，而非岩石或其他固态物质，但“超级地球”（mega-Earths）这一类别也确实存在。在太阳系中共有四颗此类行星：木星、土星、天王星和海王星。此外，已有大量系外巨行星被发现。

巨行星有时被称为气体巨行星（gas giants），但许多天文学家现在将此称谓仅用于木星和土星，并将成分不同的天王星与海王星归类为冰巨行星（ice giants）。然而，这两种名称都可能导致误解：在太阳系中，所有巨行星主要由处于临界点之上的流体组成，在此状态下气相与液相不再具有明确区分。木星和土星主要由氢与氦构成，而天王星与海王星则由水、氨和甲烷构成。

关于低质量褐矮星与高质量气体巨行星（约 \(13\,M_{\mathrm{J}}\)）之间的定义差异仍存在争议。一种观点基于行星的形成过程；另一种观点则基于行星内部物理性质。争论的一部分涉及褐矮星是否必须在其演化史中的某一阶段经历过核聚变，这一点是否应被视为定义它们的必要条件。\(^\text{[1]}\)
\subsection{术语}
“气体巨行星”（gas giant）一词由科幻作家 James Blish 于 1952 年首次提出，最初用于指代所有巨行星。然而，这个名称在某种意义上并不准确，因为在这些行星的大部分体积中，压力高到足以使物质不再以气态存在。除大气的上层区域之外，行星内部的所有物质都可能处于临界点以上的状态，在该状态下液相与气相已无法区分。\(^\text{[2]}\)因此，“流体行星”（fluid planet）会是更精确的称呼。以木星为例，其中心附近存在金属氢，但其大部分体积仍由氢、氦以及微量其他气体构成，这些物质均处在其临界点以上。\(^\text{[3]}\)所有巨行星的可观测大气层（光学深度小于 1 的区域）相对行星半径而言极为薄弱，可能只延伸至中心约百分之一的位置。因此，可观测区域确实是气态的（与火星和地球不同，它们的大气层为观察其固体地壳提供了窗口）。

这一易引起误解的术语之所以流行，是因为行星科学家通常将 rock（岩质物质）、gas（气体）与 ice（冰质物质）作为行星成分的分类速记，与这些物质在行星内部的实际相态无关。在外太阳系中，氢与氦被归类为气体；水、甲烷与氨被归类为冰；硅酸盐与金属则被归类为岩石。若从行星深部内部的角度出发，“冰”基本指含氧与含碳物质，“岩石”指含硅物质，而“气体”则指氢和氦。由于天王星与海王星在许多方面与木星和土星显著不同，一些研究者开始仅将“气体巨行星”用于后两者。在此术语体系下，部分天文学家开始将天王星与海王星称为“冰巨行星”（ice giants），以强调其内部成分中冰（以流体形态存在）的占优势地位。\(^\text{[4]}\)

另一术语“类木行星”（jovian planet）源自罗马神话的木星神 Jupiter——其属格形式为 Jovis，故称 Jovian——意在表示这一类行星都与木星相似。

对质量足以点燃氘聚变的天体（在太阳成分条件下通常需超过 13 个木星质量）称为褐矮星（brown dwarf），其质量范围介于大型巨行星与最低质量恒星之间。\(^\text{[5]}\)13 个木星质量（\(13\,M_{\mathrm{J}}\)）的界限只是经验规则，而非具有精确物理意义的临界点。质量较大的天体会燃烧其大部分氘，而质量较小者只能燃烧极少量，13 \(M_{\mathrm{J}}\) 只是两者之间的大致分界。\(^\text{[6]}\)氘的燃烧量不仅取决于质量，也取决于行星的成分，特别是氦与氘的含量。Extrasolar Planets Encyclopaedia 将天体质量上限设置为 60 \(M_{\mathrm{J}}\)，而 Exoplanet Data Explorer 则采用 24 \(M_{\mathrm{J}}\) 的上限。\(^\text{[7][8]}\)
\subsection{描述}
\begin{figure}[ht]
\centering
\includegraphics[width=8cm]{./figures/069f6f04024e47b5.png}
\caption{巨行星内部的剖面示意图。图中木星被描绘为在岩质核心之上覆盖着一层深厚的金属态氢层。} \label{fig_Giant_2}
\end{figure}
巨行星是一类质量极大的行星，并具有由氢和氦组成的厚重大气层。在形成过程中，较重元素可能会聚集成一个凝聚的“核心”。\(^\text{[9]}\)然而，这一核心可能部分或完全溶解并扩散于氢氦包层之中。\(^\text{[10][9]}\)在诸如木星和土星这类“传统”巨行星（气体巨行星）中，氢与氦构成了行星质量的大部分；而在天王星与海王星中，它们仅构成外层包壳，这两颗行星的主体则由水、氨和甲烷构成，因此日益被称为“冰巨行星”。

环绕恒星并处于极近轨道的系外巨行星是最容易被探测到的。这类行星因其极高的表面温度，被称为“热木星”（hot Jupiters）与“热海王星”（hot Neptunes）。在空间望远镜出现之前，热木星是已知最常见的系外行星类型，这是因为利用地面仪器探测它们相对容易。

巨行星通常被认为缺乏固体表面，但更精确地说，它们根本不具有真正意义上的“表面”，因为其构成气体在向外层延伸时会不断变得稀薄，最终与行星际介质难以区分。因此，能否在巨行星上“着陆”取决于其核心的大小与成分，不一定可行。
\subsection{子类型}  
\subsubsection{气体巨行星}
\begin{figure}[ht]
\centering
\includegraphics[width=8cm]{./figures/093e6de5991f9b39.png}
\caption{土星北极的极涡结构} \label{fig_Giant_3}
\end{figure}
气体巨行星主要由氢和氦组成。太阳系中的气体巨行星——木星与土星——其质量中约有 3\% 至 13\% 由较重元素构成。通常认为，气体巨行星由一层分子氢外壳包覆着液态金属氢层，而其内部可能存在一个由岩质成分构成的熔融核心。

在木星和土星的大气最外层，氢组成的气体包层中包含多层可见云层，这些云层主要由水和氨构成。金属氢层构成了两颗行星体积的绝大部分。之所以称其为“金属态”，是因为在极高压力下，氢会转变为具有电导性的物质。行星核心被认为由更重的元素组成，但其内部温度可高达约 20{,}000 K、压力极端之高，因此其物理性质仍难以准确理解。\(^\text{[11]}\)
\subsubsection{冰巨行星}
\begin{figure}[ht]
\centering
\includegraphics[width=8cm]{./figures/15aded728ec11429.png}
\caption{由哈勃望远镜拍摄的图像合成，展示了太阳系四颗巨行星在 2014 至 2024 年十年观测期间的季节性变化。} \label{fig_Giant_4}
\end{figure}
冰巨行星的内部组成与气体巨行星显著不同。太阳系中的冰巨行星——天王星和海王星——具有富含氢的大气层，从云顶一直延伸到其半径的大约 80\%（天王星）或 85\%（海王星）的位置。在此之下，它们的内部主要呈现“冰质”特征，即主要由水、甲烷和氨构成。它们也含有一定量的岩石与气体，但冰、岩石与气体的不同比例在结构上可能与纯冰相似，因此这些成分的确切比例目前尚不清楚。\(^\text{[12]}\)

天王星与海王星具有高度弥散的大气层，并含有少量甲烷，使它们呈现淡青绿色。两者的磁场均相对于自转轴有显著倾斜。

与其他巨行星不同，天王星具有极端的自转轴倾角，导致其季节变化十分强烈。两颗行星之间还存在其他细微但重要的差异。例如，尽管天王星总体质量更小，但其氢和氦的含量却多于海王星。因此，海王星密度更高，内部热量也更丰富，同时拥有更为活跃的大气层。根据 Nice 模型，海王星事实上是在比天王星更靠近太阳的位置形成的，因此应当含有更多的重元素。
\subsubsection{巨型地球 } 
“巨型地球”（mega-Earth）或“高质量固态行星”这一术语用于指代质量超过 \(10\,M_{\mathrm{E}}\) 的类地系外行星。此类行星主要由岩质物质构成，因此其密度将显著高于地球与气体巨行星。Kepler{-}10c 传统上被归类为一颗巨型地球，但后续研究表明其更可能是富含挥发物的迷你海王星（mini-Neptune）。\(^\text{[13][14]}\)一个名为“超高质量类地行星”（supermassive terrestrial planets, SMTP）的子分类曾用于指代质量超过 \(30\,M_{\mathrm{E}}\) 的巨型地球，例如 Kepler{-}145b。\(^\text{[15]}\)若干脉冲星行星（pulsar planets），如 PSR J1719{-}1438 b，被发现其质量高于木星，但半径却小于气体巨行星，因而被推测主要由结晶的钻石与氧构成。\(^\text{[16]}\)因此，它们可能是富含碳的行星级残余核心，来源于曾在与脉冲星交互期间被撕裂的伴星。\(^\text{[16]}\)然而，根据定义，这些天体更应被视为极低质量的白矮星，而非高密度的钻石行星。至于“冥界行星”（chthonian planets），\(^\text{[17]}\)如 TOI{-}849 b——即由蒸发后的气体巨行星或褐矮星剩余的岩质或金属核心构成——其质量也可能与巨型地球相当，甚至远超 \(30\,M_{\mathrm{E}}\)。\(^\text{[18][19]}\)

基于岩质行星的质量–半径关系，有研究提出：质量可达数千 \(M_{\mathrm{E}}\) 的巨大固态行星可能在大质量恒星（B 型与 O 型恒星；\(5\text{--}120\,M_{\odot}\)）周围形成，理由是此类恒星的原行星盘中可能含有足够多的重元素，而其强烈的紫外辐射与恒星风可将气体光蒸发，只留下重元素。\(^\text{[20]}\)然而，较新的研究显示，对于质量超过 \(10\,M_{\odot}\) 的恒星，其原行星盘质量与恒星质量的比例会快速下降。\(^\text{[21]}\)

根据某一模型，一种假设提出：所谓的“黑洞行星”（blanets）——本质上类似普通行星——可能围绕质量至少达百万倍太阳质量（\(10^{6}\,M_{\odot}\)）的旋转超大质量黑洞运行，并可能具有与巨型固态行星相当的质量。尽管“失控吸积”使得 blanets 从气体增长成为气体巨行星是可能的，但这一过程可能难以发生。然而，这仍然取决于围绕 blanets 的轨道区域被气体填充的速度。\(^\text{[22]}\)

